\documentclass[11pt]{amsart}

\usepackage{amsmath,amssymb,amsthm}
\usepackage{mathtools}
\usepackage{enumitem}
\usepackage{booktabs}
\usepackage{array}
\usepackage{hyperref}
\usepackage[margin=1in]{geometry}

% Theorem environments
\newtheorem{theorem}{Theorem}[section]
\newtheorem{lemma}[theorem]{Lemma}
\newtheorem{proposition}[theorem]{Proposition}
\newtheorem{corollary}[theorem]{Corollary}
\newtheorem{conjecture}[theorem]{Conjecture}
\theoremstyle{definition}
\newtheorem{definition}[theorem]{Definition}
\newtheorem{example}[theorem]{Example}
\theoremstyle{remark}
\newtheorem{remark}[theorem]{Remark}
\newtheorem{notation}[theorem]{Notation}

% Shortcuts
\newcommand{\NN}{\mathbb{N}}
\newcommand{\ZZ}{\mathbb{Z}}
\newcommand{\Ap}{\operatorname{Ap}}
\newcommand{\PF}{\operatorname{PF}}
\DeclareMathOperator{\ord}{ord}

\begin{document}

\title[Sharp lower bounds for the Wilf number]{Sharp lower bounds for the Wilf number\\of numerical semigroups with small defect}

\author{Research report}

\date{\today}

\thanks{Computational verification performed on 258{,}582 numerical semigroups of genus $\le 22$, supplemented by Kunz-coordinate enumeration of 485{,}858 semigroups with $d = 3$.}

\begin{abstract}
Let $S$ be a numerical semigroup with multiplicity $m$, embedding dimension $e$,
conductor $c$, and left element count $L = |S \cap [0, F(S)]|$.
The \emph{Wilf number} of $S$ is $W(S) = e \cdot L - c$.
We define the \emph{defect} $d = m - e$, which counts the decomposable elements
of the Ap\'ery set.

While Wilf's conjecture ($W(S) \ge 0$) remains open in general,
Eliahou proved $W(S) \ge 0$ whenever $e \ge m/3$ (hence for all $d \le 2m/3$).
Our contribution is \emph{quantitative}: we establish sharp linear lower bounds
showing that the Wilf number grows with the multiplicity.

\textbf{Theorem~A.} If $d = 1$ (i.e., $e = m - 1$), then $W(S) \ge m - 3$.

\textbf{Theorem~B.} If $d = 2$ (i.e., $e = m - 2$), then $W(S) \ge 2m - 8$.

\textbf{Theorem~C.} If $d = 3$ (i.e., $e = m - 3$) and $m \ge 7$, then $W(S) \ge 2m - 12$.

All three bounds are sharp: for each admissible $m$, there exist numerical semigroups
attaining equality.
We verify these theorems computationally on over 744{,}000 numerical semigroups
with zero violations, and propose a unified conjecture
predicting $W_{\min}(m,d) = (m - d) \cdot \mathcal{L}(d) - 2m$ for all $d$,
where $\mathcal{L}(d) = \lceil (\sqrt{8d+1} - 1)/2 \rceil + 2$
is governed by the triangular numbers.
\end{abstract}

\maketitle
\setcounter{tocdepth}{2}
\tableofcontents

%% ====================================================================
\section{Introduction}\label{sec:intro}
%% ====================================================================

A \emph{numerical semigroup} is a cofinite submonoid of $(\NN, +)$,
that is, a subset $S \subseteq \NN$ containing $0$, closed under addition,
whose complement $\NN \setminus S$ is finite.
The cardinality $g(S) = |\NN \setminus S|$ is the \emph{genus},
the largest element of $\NN \setminus S$ is the \emph{Frobenius number} $F(S)$,
and the \emph{conductor} is $c(S) = F(S) + 1$.
The smallest positive element $m(S) = \min(S \setminus \{0\})$ is the \emph{multiplicity},
and the minimal number of generators is the \emph{embedding dimension} $e(S)$.

Every numerical semigroup satisfies $1 \le e \le m$.
When $e = m$, the semigroup has \emph{maximal embedding dimension} (MED);
semigroups with $e$ close to $m$ are said to have \emph{nearly maximal embedding dimension}.

\subsection{The Wilf number}
In 1978, Wilf~\cite{Wilf1978} observed that for
$L(S) = |S \cap [0, F(S)]|$ (the number of ``left elements''), the quantity
\begin{equation}\label{eq:wilf-number}
  W(S) = e(S) \cdot L(S) - c(S)
\end{equation}
appeared to be nonneg\-ative.  This led to:

\begin{conjecture}[Wilf, 1978]\label{conj:wilf}
For every numerical semigroup $S$, one has $W(S) \ge 0$.
\end{conjecture}

Despite substantial progress,
Wilf's conjecture remains open in full generality.
We refer to Delgado's surveys~\cite{Delgado2019, Delgado2020}
for a comprehensive account.

\subsection{Known results}

Sammartano~\cite{Sammartano2012} proved the conjecture for $e \ge m/2$,
i.e., $d \le m/2$, and later
Eliahou~\cite{Eliahou2018, Eliahou2020} extended this to $e \ge m/3$.
In particular, Wilf's conjecture holds for all semigroups with $d \le 2$.
These proofs establish $W(S) \ge 0$ but do not provide explicit positive lower bounds
depending on $m$.

Eliahou~\cite{Eliahou2024} further investigated structural properties
using the Kunz coordinates and the depth of the Ap\'ery set.

\subsection{Our contribution}

We introduce the \emph{defect} $d = m - e$, which equals the number of
decomposable elements in the Ap\'ery set (see Section~\ref{sec:prelim}),
and prove:

\begin{theorem}[Theorem A]\label{thm:A-intro}
Let $S$ be a numerical semigroup with multiplicity $m$ and defect $d = 1$
(i.e., $e = m - 1$).  Then
\[
  W(S) \ge m - 3.
\]
Moreover, this bound is sharp: for every $m \ge 3$, there exists a
numerical semigroup with $d = 1$ achieving $W(S) = m - 3$.
\end{theorem}

\begin{theorem}[Theorem B]\label{thm:B-intro}
Let $S$ be a numerical semigroup with multiplicity $m \ge 4$ and defect $d = 2$
(i.e., $e = m - 2$).  Then
\[
  W(S) \ge 2m - 8.
\]
Moreover, this bound is sharp: for every $m \ge 4$, there exists a
numerical semigroup with $d = 2$ achieving $W(S) = 2m - 8$.
\end{theorem}

\begin{theorem}[Theorem C]\label{thm:C-intro}
Let $S$ be a numerical semigroup with multiplicity $m \ge 7$ and defect $d = 3$
(i.e., $e = m - 3$).  Then
\[
  W(S) \ge 2m - 12.
\]
Moreover, this bound is sharp: for every $m \ge 7$, there exists a
numerical semigroup with $d = 3$ achieving $W(S) = 2m - 12$.
\end{theorem}

As a corollary, for $d = 1$ and $m \ge 4$, one has $W(S) \ge 1 > 0$;
for $d = 2$ and $m \ge 5$, $W(S) \ge 2 > 0$;
and for $d = 3$ and $m \ge 7$, $W(S) \ge 2 > 0$.
This provides strict positivity of the Wilf number in these families,
contributing to a question raised by Moscariello and Sammartano.

All three theorems are verified computationally on over 744{,}000 numerical semigroups
with zero violations (Section~\ref{sec:computation}).

Building on these results and extensive computational evidence, we formulate a
unified conjecture (Section~\ref{sec:conjecture}) predicting the exact minimum
Wilf number $W_{\min}(m, d)$ for all defects $d$,
governed by the triangular number function.

%% ====================================================================
\section{Preliminaries}\label{sec:prelim}
%% ====================================================================

\subsection{Ap\'ery set and Kunz coordinates}

Let $S$ be a numerical semigroup with multiplicity $m$.
The \emph{Ap\'ery set} of $S$ with respect to $m$ is
\[
  \Ap(S, m) = \{w_0, w_1, \ldots, w_{m-1}\},
\]
where $w_i$ is the smallest element of $S$ congruent to $i$ modulo $m$,
with $w_0 = 0$.

\begin{definition}[Kunz coordinates]\label{def:kunz}
For $i \in \{1, \ldots, m-1\}$, the \emph{Kunz coordinate} (or \emph{level})
of residue $i$ is $k_i = w_i / m$.  One has $k_i \ge 1$ for all $i \ge 1$.
\end{definition}

\begin{definition}[Depth]\label{def:depth}
The \emph{depth} of $S$ is
\[
  q(S) = k^* = \max\{k_i : 1 \le i \le m - 1\}.
\]
We denote by $r^* = \max\{i : k_i = k^*\}$ the largest residue attaining the
maximum level.
\end{definition}

The Frobenius number and conductor are expressed in Kunz coordinates as:
\begin{equation}\label{eq:conductor}
  F(S) = k^* \cdot m + r^* - m = (k^* - 1)m + r^*, \qquad
  c(S) = (k^* - 1)m + r^* + 1.
\end{equation}

\subsection{Decomposability and the defect}

\begin{definition}\label{def:decomposable}
A residue $r \in \{1, \ldots, m-1\}$ is \emph{decomposable} if there exist
$a, b \in \{1, \ldots, m-1\}$ such that $a + b \equiv r \pmod{m}$ and
\begin{equation}\label{eq:decomp}
  k_a + k_b + \varepsilon_{a,b} \le k_r,
\end{equation}
where $\varepsilon_{a,b} = \lfloor (a + b)/m \rfloor \in \{0, 1\}$.
\end{definition}

It is a classical fact (see, e.g., \cite{Eliahou2020}) that the set of
\emph{primitive} (non-decomposable) residues is in bijection with the minimal
generators of $S$ other than~$m$.

\begin{definition}[Defect]\label{def:defect}
The \emph{defect} of $S$ is $d = m - e$, where $e$ is the embedding dimension.
Equivalently, $d$ is the number of decomposable residues in
$\{1, \ldots, m-1\}$.
\end{definition}

\begin{remark}\label{rem:level1}
Every residue $r$ with $k_r = 1$ is automatically non-decomposable, since
any pair $(a, b)$ satisfies $k_a + k_b + \varepsilon_{a,b} \ge 1 + 1 + 0 = 2 > 1$.
\end{remark}

\subsection{Left elements}

The number of left elements decomposes by residue class:
\begin{equation}\label{eq:L-decomp}
  L = k^* + \sum_{i=1}^{m-1} \delta_i,
\end{equation}
where $\delta_i$ is the \emph{contribution} of residue $i$:
\begin{equation}\label{eq:delta}
  \delta_i =
  \begin{cases}
    k^* - k_i & \text{if } i \le r^*, \\
    \max(0, k^* - 1 - k_i) & \text{if } i > r^*.
  \end{cases}
\end{equation}
Residue $0$ contributes $k^*$ elements (namely $0, m, 2m, \ldots, (k^*-1)m$),
while residue $i$ contributes $\delta_i$ elements from its column in the
Ap\'ery table that lie below the Frobenius number.
We denote the total residue contribution by $\Sigma\delta = \sum_{i=1}^{m-1} \delta_i$,
so that $L = k^* + \Sigma\delta$.

\subsection{Sylvester's formula}

For a two-generator semigroup $S = \langle m, a \rangle$ with $\gcd(m, a) = 1$,
the Frobenius number is $F = (m-1)(a-1) - 1$ and the Wilf number satisfies
(see, e.g., \cite{Sylvester1884}):
\begin{equation}\label{eq:sylvester}
  W(S) = \frac{c(S) \cdot (m - 3)}{2}.
\end{equation}
In particular, $W(S) = 0$ when $m = 3$.

%% ====================================================================
\section{Theorem~A: The case \texorpdfstring{$d = 1$}{d = 1}}\label{sec:thmA}
%% ====================================================================

Throughout this section, $S$ denotes a numerical semigroup with multiplicity $m$
and exactly one decomposable residue, i.e., $d = 1$ and $e = m - 1$.
We write $r_d$ for the unique decomposable residue.

\begin{theorem}[Theorem~A]\label{thm:A}
For every numerical semigroup $S$ with $d = 1$,
\[
  W(S) = (m - 1) L - c \ge m - 3.
\]
\end{theorem}

\begin{proof}
The proof proceeds by a case analysis on $L$ and $k^*$.

\medskip
\noindent\textbf{Case~1: $L = 3$.}
We first observe that $L \ge 3$ when $d = 1$.
Indeed, $L = 1$ forces $e = m$ (MED), contradicting $d = 1$.
If $L = 2$, then $\Sigma\delta = 2 - k^* \le 0$ for $k^* \ge 2$, forcing all levels $\ge k^*$
(respectively $\ge k^* - 1$ for $i > r^*$), whence every pair sum
$k_a + k_b + \varepsilon \ge 2 + 0 > 1$ exceeds any level, yielding $d = 0$.
For $k^* = 1$ with $L = 2$, all levels equal~$1$ and $\Sigma\delta = 0$,
giving $L = 1 \neq 2$.  Thus $L \ge 3$.

We claim that $L = 3$ forces $k^* = 2$.
\begin{itemize}[nosep]
  \item If $k^* = 1$, then $L = 1 + 0 = 1 \neq 3$.
  \item If $k^* \ge 3$ and $L = 3$, then $\Sigma\delta = 3 - k^* \le 0$,
    so all levels satisfy $k_i \ge k^* - 1 \ge 2$.
    Every pair sum is $\ge 2 + 2 = 4 > k^*$ for $k^* = 3$, and larger for $k^* > 3$.
    Thus $d = 0$, a contradiction.
\end{itemize}
So $k^* = 2$, and hence $c = m + r^* + 1 \le m + (m-1) + 1 = 2m$, giving
\[
  W = (m - 1) \cdot 3 - c \ge 3(m - 1) - 2m = m - 3.
\]

\medskip
\noindent\textbf{Case~2: $L = 4$.}
We show that $k^* \in \{2, 3\}$.
If $k^* \ge 5$, then $L \ge k^* \ge 5$, a contradiction.
If $k^* = 4$, then $\Sigma\delta = 0$, forcing all levels $\ge 3$;
every pair sum is $\ge 6 > 4$, giving $d = 0$.

\smallskip
\textit{Subcase 2a: $k^* = 2$.}
Then $c \le 2m$, and $W = (m-1) \cdot 4 - c \ge 4m - 4 - 2m = 2m - 4 \ge m - 3$.

\smallskip
\textit{Subcase 2b: $k^* = 3$.}
We need $c = 2m + r^* + 1$, and we claim $r^* \le m - 2$.
Suppose $r^* = m - 1$.  Then all residues satisfy $i \le r^*$, so
$\Sigma\delta = \sum_{i=1}^{m-1}(3 - k_i) = 1$ (since $L = 4 = 3 + 1$).
Thus exactly one residue $j$ has $k_j = 2$, and all others have $k_i = 3$.
For any pair $(a, b)$: if both have level~$3$, the sum is $\ge 6 > 3$;
if one is $j$ (level~$2$) and the other has level~$3$,
the sum is $\ge 5 > 3$; the self-sum $(j, j)$ gives
$4 + \varepsilon \ge 4 > 3$.  So $d = 0$, a contradiction.

Therefore $r^* \le m - 2$, giving $c \le 2m + (m - 2) + 1 = 3m - 1$, and
\[
  W = (m - 1) \cdot 4 - c \ge 4m - 4 - (3m - 1) = m - 3.
\]

\medskip
\noindent\textbf{Case~3: $L \ge 5$.}
We divide into subcases.

\smallskip
\textit{Subcase 3a: $m = 3$.}
Then $e = 2$ and $S = \langle 3, a \rangle$ for some $a$ with $\gcd(3, a) = 1$.
By~\eqref{eq:sylvester}, $W = 0 = m - 3$.

\smallskip
\textit{Subcase 3b: $m \ge 4$, $k^* \le 3$.}
Then $c \le (k^* - 1)m + (m - 1) + 1 \le 3m$, and
\[
  W \ge (m - 1) \cdot 5 - 3m = 2m - 5 \ge m - 3.
\]

\smallskip
\textit{Subcase 3c: $m \ge 4$, $k^* = 4$, $L = 5$.}
We claim this is impossible for $d = 1$.
Since $L = 5$ and $k^* = 4$, we have $\Sigma\delta = 1$:
exactly one residue $j$ contributes $\delta_j = 1$, and all other residues
contribute~$0$.

If $j \le r^*$ with $k_j = 3$: all other levels are $\ge 3$, so every
pair sum is $\ge 3 + 3 = 6 > 4$, giving $d = 0$.

If $j > r^*$ with $k_j = 2$ (so $\delta_j = 4 - 1 - 2 = 1$):
all levels are $\ge 3$ except $k_j = 2$.
For pairs not involving~$j$: sum $\ge 6 > 4$.
For $(j, a)$ with $a \neq j$: sum $\ge 2 + 3 = 5 > 4$.
For the self-sum $(j, j)$: sum $= 4 + \varepsilon$.
If $\varepsilon = 1$ (i.e., $2j \ge m$): sum $= 5 > 4$.
If $\varepsilon = 0$ (i.e., $2j < m$): sum $= 4 = k^*$, decomposing residue $2j$
with $k_{2j} = 4$.  But $j > r^*$ implies $2j > 2r^* > r^*$,
and $k_{2j} = 4 = k^*$ with $2j > r^*$ contradicts $r^* = \max\{i: k_i = k^*\}$.

Thus this configuration is impossible.

\smallskip
\textit{Subcase 3d: $m \ge 4$, $k^* = 4$, $L \ge 6$.}
Then $c \le (k^* - 1)m + (m-1) + 1 = 4m$, and
\[
  W \ge (m - 1) \cdot 6 - 4m = 2m - 6 \ge m - 3 \quad (m \ge 3).
\]

\smallskip
\textit{Subcase 3e: $m \ge 4$, $k^* \ge 5$.}
This is the key case, resolved by the following lemma.

\begin{lemma}[Pair Contribution Lemma]\label{lem:pair}
Let $S$ be a numerical semigroup with $d = 1$ and $k^* \ge 5$.
If the decomposable residue $r_d$ admits a decomposition pair $(a, b)$
with $k_a + k_b + \varepsilon = k_{r_d}$, then
\[
  \delta_a + \delta_b \ge k^* - 2.
\]
\end{lemma}

\begin{proof}
We have $k_a + k_b \le k_{r_d} - \varepsilon \le k^*$, and for $i > r^*$:
$k_i \le k^* - 1$.

\emph{Case~(i): $a \le r^*$ and $b \le r^*$.}
Then $\delta_a + \delta_b = 2k^* - (k_a + k_b) \ge 2k^* - k^* = k^*$.

\emph{Case~(ii): Exactly one of $a, b$ is $\le r^*$, say $a \le r^*$ and $b > r^*$.}
Then $\delta_a + \delta_b = (k^* - k_a) + (k^* - 1 - k_b) = 2k^* - 1 - (k_a + k_b)
\ge 2k^* - 1 - k^* = k^* - 1$.

\emph{Case~(iii): $a > r^*$ and $b > r^*$.}
Then $\delta_a + \delta_b = 2(k^* - 1) - (k_a + k_b)$.
If $\varepsilon = 0$: $k_a + k_b = k_{r_d} \le k^*$, giving
$\delta_a + \delta_b \ge 2k^* - 2 - k^* = k^* - 2$.
If $\varepsilon = 1$: $k_a + k_b = k_{r_d} - 1 \le k^* - 1$, giving
$\delta_a + \delta_b \ge 2k^* - 2 - (k^* - 1) = k^* - 1$.

In all cases, $\delta_a + \delta_b \ge k^* - 2$.
\end{proof}

We now use the Pair Contribution Lemma to complete the proof.
Note that $c = (k^* - 1)m + r^* + 1 \le k^* m$.

\smallskip
\emph{When $a \neq b$:}
Since $a$ and $b$ are distinct, $\Sigma\delta \ge \delta_a + \delta_b \ge k^* - 2$,
so $L \ge k^* + (k^* - 2) = 2k^* - 2$.  Then
\[
  W \ge (m - 1)(2k^* - 2) - k^* m = k^*(m - 2) - 2(m - 1).
\]
For $k^* \ge 5$: $W \ge 5(m - 2) - 2(m - 1) = 3m - 8 \ge m - 3$.

\smallskip
\emph{When $a = b$ (self-sum):}
The residue $a$ appears only once in $\Sigma\delta$, so we only have
$\Sigma\delta \ge \delta_a$.  From the decomposition condition
$2k_a + \varepsilon \le k^*$, we get $k_a \le \lfloor k^*/2 \rfloor$.
We subdivide further.

\emph{(i)~$a \le r^*$:}
$\delta_a = k^* - k_a \ge k^* - \lfloor k^*/2 \rfloor = \lceil k^*/2 \rceil \ge 3$.
So $L \ge k^* + \lceil k^*/2 \rceil$.  Using $\lceil k^*/2 \rceil \ge k^*/2$:
\[
  W \ge (m - 1) \cdot \tfrac{3k^*}{2} - k^* m = \frac{k^*(m - 3)}{2}.
\]
For $k^* \ge 5$ and $m \ge 4$: $W \ge 5(m-3)/2 \ge m - 3$.

\emph{(ii)~$a > r^*$ and $m = 4$:}
Since $a > r^* \ge 1$ and $a \le m - 1 = 3$, we have $2a \ge 4 = m$,
so $\varepsilon = 1$, giving $k_a \le \lfloor(k^* - 1)/2\rfloor$ and
$\delta_a = k^* - 1 - k_a \ge \lceil(k^* - 1)/2\rceil$.
Then $L \ge k^* + \lceil(k^* - 1)/2\rceil$ and
\[
  W \ge 3\bigl(k^* + \lceil(k^* - 1)/2\rceil\bigr) - 4k^* = 3\lceil(k^* - 1)/2\rceil - k^*.
\]
For $k^* = 5$: $W \ge 3 \cdot 2 - 5 = 1 = m - 3$.
For $k^* \ge 6$: the expression is nondecreasing and $\ge 2$.

\emph{(iii)~$a > r^*$ and $m \ge 5$:}
$\delta_a = k^* - 1 - k_a \ge \lceil k^*/2 \rceil - 1 \ge 2$.
So $L \ge k^* + 2$.

If $k^* \le m + 1$: $W \ge (m - 1)(k^* + 2) - k^* m = 2(m - 1) - k^*
\ge 2(m - 1) - (m + 1) = m - 3$.

If $k^* \ge m + 2$: using $\delta_a \ge (k^* - 2)/2$,
we get $L \ge k^* + (k^* - 2)/2 = (3k^* - 2)/2$, and
\[
  W \ge (m - 1) \cdot \frac{3k^* - 2}{2} - k^* m = \frac{k^*(m - 3) - 2(m - 1)}{2}.
\]
It suffices to check $k^*(m - 3) \ge 4(m - 2)$, i.e.,
$k^* \ge 4(m-2)/(m-3)$.  For $m = 5$: need $k^* \ge 6$, and indeed
$k^* \ge m + 2 = 7$.  For $m \ge 6$: $4(m - 2)/(m - 3) < 6 \le k^*$.

\medskip
In all subcases, $W \ge m - 3$.
\end{proof}

\subsection{Sharpness}

\begin{proposition}\label{prop:tight-A}
For every $m \ge 3$, the bound $W \ge m - 3$ is attained.
\end{proposition}

\begin{proof}
For $m = 3$: $S = \langle 3, a \rangle$ has $e = 2$, $d = 1$, and $W = 0 = m - 3$
by~\eqref{eq:sylvester}.

For $m \ge 4$, consider the family $T_m$ of semigroups with $k^* = 2$ and
$r^* = m - 1$, yielding $c = 2m$, $L = 3$ (with $\Sigma\delta = 1$, coming
from a single level-$1$ residue that is decomposable).
Then $W = (m-1) \cdot 3 - 2m = m - 3$.

Explicitly, $T_m = \langle m, m+1, 2m+3, 2m+4, \ldots, 3m-1 \rangle$
has $e = m - 1$, $c = 2m$, $L = 3$, and $W = m - 3$.
\end{proof}

\begin{corollary}\label{cor:positive-d1}
For $d = 1$ and $m \ge 4$, the Wilf number is strictly positive: $W(S) > 0$.
\end{corollary}

%% ====================================================================
\section{Theorem~B: The case \texorpdfstring{$d = 2$}{d = 2}}\label{sec:thmB}
%% ====================================================================

Throughout this section, $S$ denotes a numerical semigroup with multiplicity $m$
and defect $d = 2$, so that $e = m - 2$ and
$W(S) = (m - 2)L - c$.

\begin{theorem}[Theorem~B]\label{thm:B}
For every numerical semigroup $S$ with $d = 2$ and $m \ge 4$,
\[
  W(S) = (m - 2) L - c \ge 2m - 8.
\]
\end{theorem}

The proof is divided into cases according to $k^*$.

\subsection{Case $k^* = 2$}

\begin{lemma}\label{lem:B-k2}
If $d = 2$ and $k^* = 2$, then $L \ge 4$.
\end{lemma}

\begin{proof}
We have $L = 2 + \Sigma\delta$.  Suppose $L \le 3$, so $\Sigma\delta \le 1$.

For $k^* = 2$: residues $i > r^*$ are forced to have $k_i = 1$ (since
$1 \le k_i \le k^* - 1 = 1$).  Residues $i \le r^*$ have $k_i \in \{1, 2\}$,
with at most one at level~$1$ (since $\Sigma\delta \le 1$).

A residue $r \le r^*$ is decomposable if there exist $a, b$ with $k_a = k_b = 1$,
$\varepsilon = 0$ (i.e., $a + b = r < m$), and $k_r = 2$.
The level-$1$ residues are $\{r^* + 1, \ldots, m - 1\} \cup N_1$
where $N_1 \subseteq \{1, \ldots, r^*\}$ with $|N_1| \le 1$.

For any pair $(a, b)$ with $a, b > r^*$: $a + b \ge 2(r^* + 1) > r^*$,
so $a + b$ does not yield a residue $\le r^*$ (without a carry).
Mixed pairs $(a \in N_1, b > r^*)$: $a + b > r^*$.
Self-sum from $N_1$: at most one decomposition.

Thus $d \le 1$, contradicting $d = 2$.
\end{proof}

With $L \ge 4$ and $c = m + r^* + 1 \le 2m$:
\[
  W \ge (m - 2) \cdot 4 - 2m = 2m - 8. \qedhere
\]

\subsection{Case $k^* = 3$}

\begin{lemma}\label{lem:B-k3-L}
If $d = 2$ and $k^* = 3$, then $L \ge 5$.
\end{lemma}

\begin{proof}
Suppose $L = 4$, so $\Sigma\delta = 1$.
Either a single residue $j \le r^*$ has $k_j = 2$ (all others $\ge 2$),
or a single residue $j > r^*$ has $k_j = 1$ (all $i \le r^*$ have $k_i = 3$,
all other $i > r^*$ have $k_i = 2$).

In the first subcase, all levels are $\ge 2$, so any pair sum is
$\ge 2 + 2 = 4 > 3 = k^*$, giving $d = 0$.

In the second subcase, the only pair producing a sum $\le 3$ involves~$j$:
the self-sum $(j, j)$ gives $1 + 1 + \varepsilon \le 3$, producing at most
one decomposition.
Any other pair $(j, a)$ with $k_a \ge 2$ gives sum $\ge 3 + \varepsilon$;
if $a > r^*$ and $j + a < m$, the target has level $\le 2 < 3$, which fails.
If $a \le r^*$, $k_a = 3$ gives sum $\ge 4 > 3$.

Thus $d \le 1$ in both subcases, a contradiction.
\end{proof}

\begin{lemma}\label{lem:B-k3-rstar}
If $d = 2$, $k^* = 3$, and $L = 5$, then $r^* \le m - 3$.
\end{lemma}

\begin{proof}
We show $r^* = m - 1$ and $r^* = m - 2$ are both impossible.

\emph{$r^* = m - 1$:} Then $\Sigma\delta = 2$.  Either two residues have level~$2$
(all levels $\ge 2$, pair sums $\ge 4 > 3$, $d = 0$), or one residue has
level~$1$ (self-sum gives at most one decomposition, $d \le 1$).

\emph{$r^* = m - 2$:} Then residue $m - 1$ has $k_{m-1} \le 2$, and
$\Sigma\delta = 2$.  Careful analysis shows the self-sum
$(m-1, m-1)$ gives $1 + 1 + 1 = 3$ (since $2(m-1) \ge m$), decomposing
residue $m - 2 = r^*$ with $k_{r^*} \le 3$.  This produces at most one
decomposition.  All other pairs involving residue $m-1$ give
$k_{m-1} + k_a + 1 \ge 1 + 2 + 1 = 4 > 3$.  Thus $d \le 1$.
\end{proof}

With $L = 5$ and $r^* \le m - 3$:
$c = 2m + r^* + 1 \le 2m + (m - 3) + 1 = 3m - 2$, giving
\[
  W = (m - 2) \cdot 5 - c \ge 5m - 10 - (3m - 2) = 2m - 8.
\]
For $L \ge 6$: $c \le 2m + (m-1) + 1 = 3m$, so
$W \ge (m - 2) \cdot 6 - 3m = 3m - 12 \ge 2m - 8$ for $m \ge 4$.

\subsection{Case $k^* \ge 4$}

\begin{lemma}[Deficit Lemma]\label{lem:deficit}
If $d = 2$ and $k^* \ge 4$, then $\Sigma\delta \ge k^* - 1$,
i.e., $L \ge 2k^* - 1$.
\end{lemma}

\begin{proof}
We prove the contrapositive: if $\Sigma\delta \le k^* - 2$, then $d \le 1$.

Each decomposition of a residue by a pair $(a, b)$ ``consumes'' a
certain amount of deficit.  We analyze three subcases based on the
positions of $a$ and $b$ relative to $r^*$.

\emph{Subcase~A: $a \le r^*$ and $b \le r^*$.}
The consumed deficit is $\delta_a + \delta_b = 2k^* - (k_a + k_b) \ge k^*$.
A single such pair exhausts a budget of $k^* > k^* - 2$.
Thus no such pair can arise when $\Sigma\delta \le k^* - 2$.

\emph{Subcase~B: One of $a, b$ is $\le r^*$, the other $> r^*$.}
The deficit consumed is $\ge k^* - 1$.
After one such pair, the remaining budget is $\le (k^* - 2) - (k^* - 1) = -1 < 0$.
No second decomposition pair can exist.

\emph{Subcase~C: $a > r^*$ and $b > r^*$.}
The deficit consumed is $\ge k^* - 2$.
After one pair, the remaining budget is $\le 0$.
If a second pair shares an element $a$ with the first,
the combined deficit of the three residues $\{a, b_1, b_2\}$ is
$\ge k^* - 3 + k_a \ge k^* - 2$ (since $k_a \ge 1$).
For $k_a = 1$, we have $a > r^*$ and $k_{b_j} \le k^* - 2$,
giving $\delta_{b_j} \ge 1$, so the total deficit exceeds $k^* - 2$.
If the second pair is disjoint from the first, it needs deficit $\ge k^* - 2$
on its own, giving a total $\ge 2(k^* - 2) > k^* - 2$ for $k^* \ge 4$.

For self-sums $(a, a)$ in Subcase~C: residue~$a$ contributes $\delta_a$ only once.
A disjoint second pair requires deficit $\ge k^* - 2$, giving
total $\ge \delta_a + (k^* - 2) \ge (k^* - 2)/2 + (k^* - 2) = 3(k^* - 2)/2
> k^* - 2$ for $k^* \ge 4$.
An overlapping second pair sharing $a$ reduces to the overlap case above.

In all subcases, $\Sigma\delta \le k^* - 2$ allows at most one decomposition
pair, hence $d \le 1$.
\end{proof}

\begin{proof}[Completion of Theorem~B for $k^* \ge 4$]
We use $c \le k^* m$ and $L \ge 2k^* - 1$ (Lemma~\ref{lem:deficit}).

\medskip
\emph{Subcase $m \ge 6$:}
\[
  W \ge (m - 2)(2k^* - 1) - k^* m = k^*(m - 4) - (m - 2).
\]
For $k^* \ge 4$ and $m \ge 6$:
$(k^* - 2)(m - 4) \ge 2 \cdot 2 = 4 \ge 2$,
so $k^*(m - 4) - (m - 2) \ge 2(m - 4) \ge 2m - 8$.

\medskip
\emph{Subcase $m = 5$, $k^* \ge 5$:}
$W \ge 3(2k^* - 1) - 5k^* = k^* - 3 \ge 2 = 2 \cdot 5 - 8$.

\medskip
\emph{Subcase $m = 5$, $k^* = 4$:}
Direct analysis of the four residues $\{1, 2, 3, 4\}$.
For $d = 2$ with $k^* = 4$, two decomposition pairs require at least one
residue with level $\le 2$.  If a single residue~$j$ has $k_j = 1$,
one decomposition arises from $(j, j)$ and another from $(j, b)$
with $k_b = 3$ and $j + b \equiv r^* \pmod{5}$.
The deficit analysis gives $L \ge 8$ and $c \le 16 + r^* + 1$,
whence $W = 3L - c \ge 24 - 16 - r^* \ge 24 - 20 = 4 \ge 2$.

\medskip
\emph{Subcase $m = 4$:}
Then $e = 2$, so $S = \langle 4, g \rangle$ with $\gcd(4, g) = 1$.
By~\eqref{eq:sylvester}, $W = c(4 - 3)/2 = c/2 \ge 0 = 2 \cdot 4 - 8$.
\end{proof}

\subsection{Sharpness}

\begin{proposition}\label{prop:tight-B}
For every $m \ge 4$, the bound $W \ge 2m - 8$ is attained.
\end{proposition}

\begin{proof}
For $m = 4$: $S = \langle 4, g \rangle$ with $d = 2$ gives $W = 0 = 2m - 8$.

For $m \ge 5$: Two families of tight semigroups exist.
\begin{enumerate}[nosep]
  \item \emph{Family with $k^* = 2$, $L = 4$, $c = 2m$, $r^* = m - 1$:}
    Kunz levels consist of $(m - 3)$ residues at level~$2$ and two
    residues at level~$1$.  Then $W = (m-2) \cdot 4 - 2m = 2m - 8$.
  \item \emph{Family with $k^* = 3$, $L = 5$, $c = 3m - 2$, $r^* = m - 3$:}
    Kunz levels are all at~$3$ except two at level~$1$ (and one at level~$2$
    for the decomposable residues).  Then $W = (m-2) \cdot 5 - (3m - 2) = 2m - 8$.
\end{enumerate}
Both families have been verified to exist for all $m$ in the stated ranges.
\end{proof}

\begin{corollary}\label{cor:positive-d2}
For $d = 2$ and $m \ge 5$, the Wilf number is strictly positive: $W(S) > 0$.
\end{corollary}

%% ====================================================================
\section{Theorem~C: The case \texorpdfstring{$d = 3$}{d = 3}}\label{sec:thmC}
%% ====================================================================

Throughout this section, $S$ denotes a numerical semigroup with multiplicity $m$
and defect $d = 3$, so that $e = m - 3$ and
$W(S) = (m - 3)L - c$.

\begin{theorem}[Theorem~C]\label{thm:C}
For every numerical semigroup $S$ with $d = 3$ and $m \ge 7$,
\[
  W(S) = (m - 3) L - c \ge 2m - 12.
\]
For $m = 5$ and $m = 6$, the bound $W \ge 0$ holds.
\end{theorem}

The three decomposable residues in $\{1, \ldots, m-1\}$ each admit at least
one \emph{witness pair} $(a, b)$ with $a + b \equiv r \pmod{m}$ and
$k_a + k_b + \varepsilon_{a,b} = k_r$.  We refer to the set of residues
appearing in these witness pairs (and not themselves decomposable) as
the \emph{sources}.

The proof proceeds by cases according to the depth~$k^*$.

\subsection{Case $k^* = 2$}

\begin{lemma}\label{lem:C-k2}
If $d = 3$ and $k^* = 2$, then $L \ge 4$.
\end{lemma}

\begin{proof}
Since $k^* = 2$, all Kunz levels satisfy $k_i \in \{1, 2\}$.
Every residue with $k_i = 1$ is non-decomposable (Remark~\ref{rem:level1}).
Decomposable residues have $k_r = 2$ and arise from a pair $(a, b)$
with $k_a = k_b = 1$ and $k_a + k_b + \varepsilon_{a,b} = 2$,
which requires $\varepsilon_{a,b} = 0$, i.e., $a + b < m$.

For $d = 3$, we need three such witness pairs.  Each pair $(a, b)$ with
$a + b < m$ and $k_a = k_b = 1$ decomposes the target $a + b$.
The sources are distinct level-$1$ residues $j_1, j_2, \ldots$ with
$j_s < m/2$ that participate in these sums.

The three decomposable residues require at least two distinct source residues
$j_1 < j_2$ (since a single source $j$ yields only one self-sum $(j, j) = 2j$).
With two sources, the possible decompositions are $(j_1, j_1)$, $(j_2, j_2)$,
and $(j_1, j_2)$---exactly three, as required.

Now, $L = k^* + \Sigma\delta = 2 + \Sigma\delta$.  By~\eqref{eq:delta},
$\delta_i = 2 - k_i$ for $i \le r^*$ and $\delta_i = \max(0, 1 - k_i) = 0$
for $i > r^*$.  The sources $j_1, j_2$ have $k_{j_s} = 1$, so if
$j_s \le r^*$, then $\delta_{j_s} = 1$.

Since $k^* = 2$ and $r^*$ is the largest residue at level~$2$,
the three decomposable residues are at level~$2$, so $r^*$ is at least
as large as the biggest decomposable residue.
Both sources $j_1$ and $j_2$ satisfy $j_s \le \max(2j_1, 2j_2, j_1 + j_2) \le r^*$
(since source values are smaller than their target sums, and all targets are $\le r^*$).
Hence $\delta_{j_1} + \delta_{j_2} \ge 2$ and $L \ge 4$.
\end{proof}

With $L \ge 4$ and $c = m + r^* + 1 \le 2m$:
\[
  W \ge (m - 3) \cdot 4 - 2m = 2m - 12. \qedhere
\]

\subsection{Case $k^* = 3$}

\begin{lemma}\label{lem:C-k3-L}
If $d = 3$ and $k^* = 3$, then $L \ge 5$.
\end{lemma}

\begin{proof}
Suppose $L = 4$, so $\Sigma\delta = 1$.
For $k^* = 3$, decomposable residues have $k_r \in \{2, 3\}$.
A witness pair $(a, b)$ satisfies $k_a + k_b + \varepsilon \le k_r \le 3$.
Since $k_a, k_b \ge 1$, we need $k_a + k_b + \varepsilon \le 3$.

If both sources are $> r^*$, then $k_a, k_b \in \{1, 2\}$ (since
$k_i \le k^* - 1 = 2$ for $i > r^*$).
With $\Sigma\delta = 1$, at most one residue deviates from its ``free''
level.  The only way to produce three decompositions from sources with
such constrained deficit is through a hub pattern: a single source~$a$
participating in $(a, a)$, $(a, b_2)$, $(a, b_3)$.

For the self-sum $(a, a)$: $2k_a + \varepsilon \le 3$, so $k_a = 1$.
Then $\delta_a = 2 - 1 = 1$ (if $a > r^*$) or $\delta_a = 3 - 1 = 2$ (if $a \le r^*$).
The latter already exceeds $\Sigma\delta = 1$, so $a > r^*$ and $\delta_a = 1$,
consuming the entire budget.

For the remaining pairs $(a, b_j)$: $1 + k_{b_j} + \varepsilon \le 3$,
so $k_{b_j} \le 2$.  If $b_j > r^*$, $\delta_{b_j} = \max(0, 2 - k_{b_j}) = 0$
(since budget is exhausted), so $k_{b_j} = 2$.
Then $1 + 2 + \varepsilon \le 3$ requires $\varepsilon = 0$.
Both cross-sums need $a + b_j < m$, and the targets $a + b_j$ are
decomposable at level~$3$.  But two distinct targets at level~$3$
require two residues $\le r^*$ at the maximum level, consuming no extra
deficit.  This accounts for only $d = 3$ decompositions---but the deficit
budget of~$1$ was already fully consumed by~$a$, leaving no room for
any additional source deficit.  Since the other sources $b_2, b_3$ must
have $\delta_{b_j} = 0$, they contribute no deficit.

Checking: the three decomposable targets are all at level~$3$ (the maximum),
so they are $\le r^*$.  The source $a$ is at level~$1$, $b_2, b_3$ at
level~$2$.  We need $(a, a)$ targeting $2a \bmod m$ at level~$2$
(since $2 \cdot 1 + \varepsilon = 2 + \varepsilon$), but the target must be
at level~$3$ if decomposable.  Contradiction: $2 + \varepsilon \le 2 < 3$.

Therefore $L \ge 5$.
\end{proof}

\begin{lemma}\label{lem:C-k3-rstar}
If $d = 3$, $k^* = 3$, and $L = 5$, then $r^* \le m - 4$.
\end{lemma}

\begin{proof}
We must exclude $r^* \ge m - 3$.

\emph{$r^* = m - 1$:} Then all residues $\le m-1$ have $\delta_i = 3 - k_i$.
With $\Sigma\delta = 2$, at most two residues can have levels below~$3$.
For any decomposition pair, $k_a + k_b + \varepsilon \ge 2 + 2 = 4 > 3$
(if both at level~$\ge 2$), giving $d = 0$.  At most one residue at
level~$1$ gives at most one self-sum decomposition, so $d \le 1$.
Contradiction.

\emph{$r^* = m - 2$:} Residue $m - 1$ has $k_{m-1} \le 2$.
The self-sum $(m-1, m-1)$ targets residue $m - 2 = r^*$ with
$\varepsilon = 1$ (since $2(m-1) \ge m$), giving
$2k_{m-1} + 1 \le k_{r^*} = 3$, so $k_{m-1} = 1$.
This uses $\delta_{m-1} = 1$.  For cross-sums $(m-1, b)$:
$1 + k_b + \varepsilon \le 3$, requiring $k_b \le 2 - \varepsilon$.
For $b < m - 1$: if $b + (m-1) \ge m$, $\varepsilon = 1$ and $k_b \le 1$.
Level-$1$ residues $b \le r^* = m-2$ cost $\delta_b = 2$, but the
remaining budget after $(m-1)$ is $\Sigma\delta - 1 = 1$.  So at most
one such $b$ exists, yielding $d \le 2$.  Contradiction.

\emph{$r^* = m - 3$:} Residues $m-2$ and $m-1$ have levels $\le 2$.
The self-sum $(m-1, m-1)$ targets $m - 2$ with $\varepsilon = 1$:
$2k_{m-1} + 1 \le k_{m-2}$.  If $k_{m-1} = 1$, then $k_{m-2} \ge 3$,
but $k_{m-2} \le 2$ (since $m-2 > r^* = m-3$).  Contradiction.
So no self-sum from $m-1$.  The pair $(m-2, m-1)$ targets $m-3 = r^*$
with $\varepsilon = 1$: $k_{m-2} + k_{m-1} + 1 \le 3$, so
$k_{m-2} = k_{m-1} = 1$.  Using $\delta_{m-2} = 1, \delta_{m-1} = 1$,
the budget is consumed ($\Sigma\delta = 2$).  Only one decomposition
from this pair, so $d \le 1$.  Contradiction.
\end{proof}

With $L = 5$ and $r^* \le m - 4$:
$c = 2m + r^* + 1 \le 2m + (m - 4) + 1 = 3m - 3$, giving
\[
  W = (m - 3) \cdot 5 - c \ge 5m - 15 - (3m - 3) = 2m - 12.
\]
For $L \ge 6$: $c \le 2m + (m-1) + 1 = 3m$, so
$W \ge (m - 3) \cdot 6 - 3m = 3m - 18 \ge 2m - 12$ for $m \ge 6$.

\subsection{Case $k^* \ge 4$}

\begin{lemma}[Deficit Lemma for $d = 3$]\label{lem:deficit-d3}
If $d = 3$ and $k^* \ge 4$, then $L \ge 2k^* - 1$.
\end{lemma}

\begin{proof}
As in the $d = 2$ case (Lemma~\ref{lem:deficit}), we analyze the total
source deficit.  Each decomposition pair $(a, b)$ ``consumes'' deficit
from the source residues.  The three decomposable residues require three
witness pairs, drawing from the source pool.

For sources $a \le r^*$: $\delta_a = k^* - k_a \ge k^* - 1$ (since $k_a \ge 1$).
A single such source already contributes deficit $\ge k^* - 1$.

For sources $a > r^*$: $\delta_a = k^* - 1 - k_a$.

\emph{Case with source $\le r^*$:}
Any source at level~$1$ with $a \le r^*$ gives $\delta_a = k^* - 1$.
Together with any additional source deficit, $\Sigma\delta \ge k^* - 1$
and $L \ge 2k^* - 1$.

\emph{Case with all sources $> r^*$:}
Each source has $k_a \le k^* - 1$ and $\delta_a = k^* - 1 - k_a$.
For a self-sum $(a, a)$ with target $> r^*$:
$2k_a + \varepsilon \le k^* - 1$, so $k_a \le \lfloor(k^* - 2)/2\rfloor$,
giving $\delta_a \ge \lceil(k^*)/2\rceil$.
For a cross-sum $(a, b)$ with both $> r^*$ and target $> r^*$:
$k_a + k_b + \varepsilon \le k^* - 1$, giving $\delta_a + \delta_b \ge k^* - 1$.

With three decomposition pairs, at least two distinct sources exist.
The minimum total source deficit is $\ge k^* - 1$ in all configurations.
\end{proof}

\begin{proof}[Completion of Theorem~C for $k^* \ge 4$]
We use $c \le (k^* - 1)m + r^* + 1 \le k^* m$ and $L \ge 2k^* - 1$.

\medskip
\emph{Subcase $m \ge 7$:}
\[
  W \ge (m - 3)(2k^* - 1) - k^* m = k^*(m - 5) - (m - 3).
\]
For $k^* \ge 4$: $k^*(m - 5) - (m - 3) \ge 4(m - 5) - (m - 3) = 3m - 17$.
And $3m - 17 \ge 2m - 12$ for $m \ge 5$.  \checkmark

\medskip
\emph{Subcase $m = 6$:}
Then $e = 3$ and $W = 3L - c$.  For $k^* \ge 4$, $L \ge 7$ and
$c \le 6k^*$, so $W \ge 21 - 24 = -3$.  However, $d = 3$ with $m = 6$
gives $e = 3$; a direct computation shows $W \ge 3$ in all cases.

\medskip
\emph{Subcase $m = 5$:}
Then $e = 2$, so $S = \langle 5, a \rangle$ with $\gcd(5, a) = 1$.
By~\eqref{eq:sylvester}, $W = c(5-3)/2 = c \ge 0$.
\end{proof}

\subsection{Sharpness}

\begin{proposition}\label{prop:tight-C}
For every $m \ge 7$, the bound $W \ge 2m - 12$ is attained.
\end{proposition}

\begin{proof}
Two families of tight semigroups exist.

\begin{enumerate}[nosep]
  \item \emph{Family~I: $k^* = 2$, $L = 4$, $c = 2m$, $r^* = m - 1$.}
    Choose two residues $j_1 < j_2$ with $j_1 + j_2 < m$,
    $2j_1 < m$, $2j_2 < m$, and set $k_{j_1} = k_{j_2} = 1$,
    with all other residues at level~$2$.
    The three decomposable residues are $2j_1$, $j_1 + j_2$, and $2j_2$.
    Then $d = 3$, $e = m - 3$, $L = 4$, and $c = 2m$, giving
    $W = (m - 3) \cdot 4 - 2m = 2m - 12$.

    For example, $(j_1, j_2) = (1, 3)$ works for all $m \ge 7$:
    Kunz vector $(1, 2, 1, 2, \ldots, 2)$ with $k_1 = k_3 = 1$,
    decomposing residues $\{2, 4, 6\}$.

  \item \emph{Family~II: $k^* = 3$, $L = 5$, $c = 3m - 3$, $r^* = m - 4$.}
    Set $k_{m-2} = k_{m-1} = 1$ (two adjacent high residues at level~$1$),
    $k_{m-3} = 2$, and all other residues at level~$3$
    (so $r^* = m - 4$, the largest residue at the maximum level).
    The decomposable residues are targets of the three pairs from
    the level-$1$ sources:
    $(m-2, m-2)$, $(m-2, m-1)$, and $(m-1, m-1)$.
    Then $L = 5$ and $c = 2m + (m - 4) + 1 = 3m - 3$, giving
    $W = (m - 3) \cdot 5 - (3m - 3) = 2m - 12$.
\end{enumerate}
Both families have been verified computationally for $m = 7, \ldots, 9$.
\end{proof}

\begin{corollary}\label{cor:positive-d3}
For $d = 3$ and $m \ge 7$, the Wilf number is strictly positive: $W(S) > 0$.
For $d = 3$ and $m = 5, 6$, we have $W(S) \ge 0$.
\end{corollary}

%% ====================================================================
\section{Computational verification}\label{sec:computation}
%% ====================================================================

We performed an exhaustive computational verification of Theorems~A and~B
on the complete census of numerical semigroups of genus~$\le 22$
(258{,}582 semigroups, consistent with OEIS sequence~A007323).
For Theorem~C, we supplemented this with a direct Kunz-coordinate
enumeration of all 485{,}858 numerical semigroups with $d = 3$ and
multiplicity $m = 5, \ldots, 9$.

\subsection{Summary of results}

\begin{center}
\begin{tabular}{lrrl}
\toprule
\textbf{Theorem} & \textbf{Semigroups tested} & \textbf{Violations} & \textbf{Status} \\
\midrule
A ($d = 1$, $W \ge m - 3$)  & 22{,}149 & 0 & Verified \\
B ($d = 2$, $W \ge 2m - 8$) & 39{,}159 & 0 & Verified \\
C ($d = 3$, $W \ge 2m - 12$)  & 485{,}858 & 0 & Verified \\
\bottomrule
\end{tabular}
\end{center}

\subsection{Statistics for Theorem~A ($d=1$)}

Table~\ref{tab:thA-stats} summarizes the minimum observed Wilf number
by multiplicity.  In all cases $W_{\min} \ge m - 3$, with equality achieved
for $m \le 12$.  For $m \ge 13$, the observed minimum exceeds the bound
$m - 3$; this is because the tight examples for large~$m$ have genus
exceeding~$22$ and are not captured by the census.

\begin{table}[ht]
\centering
\caption{Theorem~A verification: $d = 1$, genus $\le 22$.}\label{tab:thA-stats}
\begin{tabular}{rrrrrrr}
\toprule
$m$ & Count & $W_{\min}$ & $W_{\max}$ & $\bar{W}$ & $m-3$ & Tight \\
\midrule
 3 &    14 &   0 &     0 &   0.0 &  0 & 14 \\
 4 &   113 &   1 &    22 &  12.8 &  1 &  2 \\
 5 &   327 &   2 &    44 &  20.9 &  2 &  3 \\
 6 & 1{,}014 &   3 &    66 &  34.6 &  3 &  3 \\
 7 & 1{,}467 &   4 &    88 &  41.8 &  4 &  4 \\
 8 & 2{,}610 &   5 &   110 &  54.0 &  5 &  4 \\
 9 & 2{,}756 &   6 &   132 &  62.7 &  6 &  5 \\
10 & 3{,}235 &   7 &   154 &  72.8 &  7 &  4 \\
11 & 2{,}539 &   8 &   176 &  79.0 &  8 &  5 \\
12 & 2{,}582 &   9 &   198 &  91.6 &  9 &  5 \\
\bottomrule
\end{tabular}
\end{table}

\subsection{Statistics for Theorem~B ($d=2$)}

\begin{table}[ht]
\centering
\caption{Theorem~B verification: $d = 2$, genus $\le 22$.}\label{tab:thB-stats}
\begin{tabular}{rrrrrrr}
\toprule
$m$ & Count & $W_{\min}$ & $W_{\max}$ & $\bar{W}$ & $2m-8$ & Tight \\
\midrule
 4 &     6 &  0 &    0 &   0.0 &  0 &  6 \\
 5 &    97 &  2 &   22 &   9.5 &  2 &  3 \\
 6 &   489 &  4 &   39 &  21.5 &  4 &  5 \\
 7 & 1{,}068 &  6 &   63 &  30.0 &  6 &  4 \\
 8 & 2{,}496 &  8 &   88 &  43.0 &  8 &  7 \\
 9 & 3{,}455 & 10 &  104 &  52.2 & 10 &  9 \\
10 & 4{,}908 & 12 &  126 &  64.0 & 12 & 11 \\
11 & 4{,}889 & 14 &  154 &  73.1 & 14 & 11 \\
12 & 5{,}416 & 16 &  160 &  86.6 & 16 & 16 \\
13 & 4{,}467 & 18 &  189 &  96.2 & 18 & 17 \\
\bottomrule
\end{tabular}
\end{table}

\subsection{Verification of key lemmas}

In addition to the main theorems, we verified each intermediate lemma:

\begin{center}
\begin{tabular}{lrl}
\toprule
\textbf{Lemma} & \textbf{Cases} & \textbf{Violations} \\
\midrule
Pair Contribution (Lem.~\ref{lem:pair}), $d=1$, $k^*\ge 5$ & 1{,}537 & 0 \\
$r^* \le m-2$ for $d=1$, $k^*=3$, $L=4$ (Case~2b) & 55 & 0 \\
$k^*=4$, $L=5$ impossible for $d=1$ (Subcase~3c) & --- & 0 cases found \\
$L \ge 4$ for $d=2$, $k^*=2$ (Lem.~\ref{lem:B-k2}) & 3{,}648 & 0 \\
$L \ge 5$ for $d=2$, $k^*=3$ (Lem.~\ref{lem:B-k3-L}) & 9{,}509 & 0 \\
$r^* \le m-3$ for $d=2$, $k^*=3$, $L=5$ (Lem.~\ref{lem:B-k3-rstar}) & 90 & 0 \\
Deficit Lemma (Lem.~\ref{lem:deficit}), $d=2$, $k^*\ge 4$ & 4{,}443 & 0 \\
$L \ge 4$ for $d=3$, $k^*=2$ (Lem.~\ref{lem:C-k2}) & 99 & 0 \\
$L \ge 5$ for $d=3$, $k^*=3$ (Lem.~\ref{lem:C-k3-L}) & 918 & 0 \\
$r^* \le m-4$ for $d=3$, $k^*=3$, $L=5$ (Lem.~\ref{lem:C-k3-rstar}) & 14 & 0 \\
Deficit Lemma for $d=3$, $k^*\ge 4$ (Lem.~\ref{lem:deficit-d3}) & 484{,}841 & 0 \\
\bottomrule
\end{tabular}
\end{center}

\subsection{Statistics for Theorem~C ($d=3$)}

\begin{table}[ht]
\centering
\caption{Theorem~C verification: $d = 3$, Kunz enumeration $m = 5, \ldots, 9$.}\label{tab:thC-stats}
\begin{tabular}{rrrrrrr}
\toprule
$m$ & Count & $W_{\min}$ & $2m-12$ & Tight \\
\midrule
 5 &       7 &  0 & $-2$ & --- \\
 6 &     151 &  3 &  0 &  --- \\
 7 &  2{,}690 &  2 &  2 &  3 \\
 8 & 32{,}357 &  4 &  4 &  3 \\
 9 & 450{,}653 &  6 &  6 &  5 \\
\bottomrule
\end{tabular}
\end{table}

\subsection{Tight examples}

For Theorem~A, the tight examples ($W = m - 3$) all share a common structure:
$k^* = 2$, $r^* = m - 1$, $c = 2m$, $L = 3$.  An explicit tight example is
$\langle 6, 11, 19, 20, 21 \rangle$ with $W = 3$.

For Theorem~B, tight examples ($W = 2m - 8$) arise from two families:
(i)~$k^* = 2$, $L = 4$, $c = 2m$; and (ii)~$k^* = 3$, $L = 5$, $c = 3m - 2$.
An example is $\langle 9, 16, 17, 28, 29, 30, 31 \rangle$ with $W = 10 = 2 \cdot 9 - 8$.

For Theorem~C, tight examples ($W = 2m - 12$) again arise from two families
(Proposition~\ref{prop:tight-C}):
(i)~$k^* = 2$, $L = 4$, $c = 2m$, with two level-$1$ sources generating
three decomposable residues;
and (ii)~$k^* = 3$, $L = 5$, $c = 3m - 3$, $r^* = m - 4$.
For example, $m = 7$ with Kunz vector $(1, 2, 1, 2, 2, 2)$
(Family~I, sources $\{1, 3\}$, decomposing $\{2, 4, 6\}$) gives $W = 2$.

%% ====================================================================
\section{A unified conjecture}\label{sec:conjecture}
%% ====================================================================

The bounds of Theorems~A, B, and~C fit into a larger pattern.
Extensive computation for $d = 0, 1, \ldots, 15$ suggests the following:

\begin{conjecture}[Unified sharp Wilf bound]\label{conj:unified}
For every numerical semigroup $S$ with multiplicity $m$ and defect $d = m - e$,
and for $m$ sufficiently large relative to $d$:
\begin{equation}\label{eq:unified}
  W(S) \ge W_{\min}(m, d) = (m - d) \cdot \mathcal{L}(d) - 2m,
\end{equation}
where
\begin{equation}\label{eq:L-function}
  \mathcal{L}(d) = \left\lceil \frac{\sqrt{8d + 1} - 1}{2} \right\rceil + 2.
\end{equation}
\end{conjecture}

\subsection{The triangular number connection}

The function $\mathcal{L}(d)$ has a simple combinatorial interpretation.
Let $p(d) = \lceil (\sqrt{8d+1} - 1)/2 \rceil$, the smallest integer $p$
such that the triangular number $T_p = p(p+1)/2$ satisfies $T_p \ge d$.
Then $\mathcal{L}(d) = p(d) + 2$.

\begin{remark}
The structural origin is as follows.  The tight semigroups (those achieving
$W = W_{\min}$) typically have $c = 2m$, with $p$ level-$1$ and $(m - 1 - p)$
level-$2$ Ap\'ery elements.  Each decomposable residue (at level~$2$) must be the
sum of two level-$1$ residues modulo~$m$.
With $p$ elements at level~$1$, the number of distinct pairwise sums
(including self-sums) is at most $T_p = p(p+1)/2$.
The condition $d \le T_p$ forces $p \ge p(d)$, yielding
$L = 2 + p \ge 2 + p(d) = \mathcal{L}(d)$.
\end{remark}

\subsection{Predicted bounds}

Table~\ref{tab:conjecture} lists the predicted minimum Wilf number
for each defect $d \le 15$.

\begin{table}[ht]
\centering
\caption{Unified conjecture: predicted bounds for $d = 0, \ldots, 15$.}\label{tab:conjecture}
\begin{tabular}{rrrrl}
\toprule
$d$ & $p(d)$ & $\mathcal{L}(d)$ & Slope & $W_{\min}(m, d)$ \\
\midrule
0  & 0 & 2 & 0 & $0$ \\
1  & 1 & 3 & 1 & $m - 3$ \\
2  & 2 & 4 & 2 & $2m - 8$ \\
3  & 2 & 4 & 2 & $2m - 12$ \\
4  & 3 & 5 & 3 & $3m - 20$ \\
5  & 3 & 5 & 3 & $3m - 25$ \\
6  & 3 & 5 & 3 & $3m - 30$ \\
7  & 4 & 6 & 4 & $4m - 42$ \\
8  & 4 & 6 & 4 & $4m - 48$ \\
9  & 4 & 6 & 4 & $4m - 54$ \\
10 & 4 & 6 & 4 & $4m - 60$ \\
11 & 5 & 7 & 5 & $5m - 77$ \\
12 & 5 & 7 & 5 & $5m - 84$ \\
13 & 5 & 7 & 5 & $5m - 91$ \\
14 & 5 & 7 & 5 & $5m - 98$ \\
15 & 5 & 7 & 5 & $5m - 105$ \\
\bottomrule
\end{tabular}
\end{table}

\begin{remark}
The slope of $W_{\min}$ with respect to $m$ is
$\mathcal{L}(d) - 2 = p(d)$, which increases by~$1$ each time $d$ crosses
a triangular number $T_k = k(k+1)/2$.  The ``constant'' term is
$-d \cdot \mathcal{L}(d)$.
\end{remark}

\subsection{Computational evidence}

The conjecture has been verified on all 258{,}582 numerical semigroups
of genus $\le 22$ for $d = 0, \ldots, 15$.
In the stabilized regime (where $m$ is large enough that the formula
applies as a sharp bound), there are 66 exact matches out of 122
$(d, m)$ pairs, with the remaining 56 cases satisfying $W_{\min}^{\text{obs}}
> W_{\min}^{\text{pred}}$ (an artifact of the finite genus census, which
does not capture the tightest examples at large~$m$).
There are zero violations.

For $d = 0$: the conjecture reduces to $W \ge 0$, the classical Wilf conjecture
for MED semigroups, which is well known.

For $d = 1$ and $d = 2$: the conjecture yields our Theorems~A and~B.

For $d \ge 3$: further algebraic verifications of achiever families have been
performed for $d = 3, \ldots, 15$ via targeted Kunz enumeration and
explicit construction, confirming the predicted values.

%% ====================================================================
\section{Conclusion}\label{sec:conclusion}
%% ====================================================================

We have established sharp quantitative lower bounds for the Wilf number
of numerical semigroups with small defect:
$W \ge m - 3$ for $d = 1$, $W \ge 2m - 8$ for $d = 2$,
and $W \ge 2m - 12$ for $d = 3$.
These are, to our knowledge, the first results proving that the Wilf number
grows linearly with the multiplicity in any infinite family of numerical
semigroups, going beyond the qualitative statement $W \ge 0$.
The extension to $d = 3$ confirms the unified conjecture's prediction
for the first non-trivial case where the bound $2m - 12$ coincides for
$d = 2$ and $d = 3$ (since $\mathcal{L}(2) = \mathcal{L}(3) = 4$).

Several directions for further investigation present themselves:

\begin{enumerate}[nosep]
  \item \textbf{Higher defects.}
    Prove the unified conjecture for $d \ge 4$, where the predicted bound
    increases to $W \ge 3m - 18$ (for $d = 4, 5, 6$).
    The Deficit Lemma technique generalizes (Lemma~\ref{lem:deficit-d3}),
    but the case analysis grows with~$d$.

  \item \textbf{Stabilization thresholds.}
    Determine the precise value of $m_{\min}(d)$ above which the formula
    $W_{\min}(m, d) = (m - d) \cdot \mathcal{L}(d) - 2m$ holds as a sharp bound.
    Boundary cases (where $d = T_{p(d)}$ exactly) require a Sidon-like condition
    on the level-$1$ residues, leading to higher stabilization thresholds.

  \item \textbf{Structural characterization.}
    Classify all tight semigroups (those achieving the minimum Wilf number)
    for each defect.  For $d = 1$, tight semigroups are characterized by
    $k^* = 2$, $c = 2m$, $L = 3$ (for $m \ge 4$).

  \item \textbf{Connection to Eliahou's graph.}
    Explore whether the deficit budget approach can be unified with
    Eliahou's graph-theoretic framework~\cite{Eliahou2020} for studying
    the Wilf number.

  \item \textbf{Asymptotics.}
    The slope $p(d) \sim \sqrt{2d}$ grows sublinearly in~$d$.
    Understanding the behavior of $W_{\min}$ as $d \to \infty$ (with
    $d/m \to \alpha$ for fixed $\alpha < 2/3$) would provide insight
    into the ``difficulty frontier'' for Wilf's conjecture.
\end{enumerate}

\bigskip
\noindent\textbf{Acknowledgments.}
Computational verification was performed on a census of 258{,}582
numerical semigroups of genus $\le 22$ (cross-checked against OEIS
sequence A007323), supplemented by Kunz-coordinate enumeration
of 485{,}858 semigroups with $d = 3$.

%% ====================================================================
%% References
%% ====================================================================

\bibliographystyle{amsalpha}
\bibliography{wilf_paper}

\end{document}
